\scanpage{71}
\mspage{70}
% E: deletion|upper left|cas au moins c'est dû au fait que
\deleted{cas au moins, c'est dû au fait que je n'ai pas tenu compte des arcs de spécialisation exceptionnels (vers des p.r. supersingulières). Le \unsure{mauvais} $\unsure{\gamma|\tau}$ \unsure{à deux faces|à deux visites}, vient qu'on se permet de se spécialiser en une pos. rel. supersingulière, \unsure{ou} qu'on veut l'éviter, en passant directement à la position voisine opposée.}
% E: text|left middle|Je viens de vérifier sauf erreur
Je viens de vérifier, sauf erreur (j'ai fait pas mal d'erreurs encore dans le cours de la vérification!) les trois procédés $A_1,A_2,B_{\unsure{1|3}}$ donnent lieu aux bons circuits autour d'une position supersingulière, il n'y a que $B_2$ qui \unsure{déconne|déraille} complètement (or il semblerait que c'était sur $B_2$ que je m'étais arrêté pendant une quinzaine, Dieu sait pourquoi!)
% E: text|bottom left|Le plus intéressant après A2
Le plus intéressant, après $A_2$, reste bien sûr $A_1$, mais je viens encore d'essayer de comprendre les circuits auxquels il donne naissance, sans y réussir.
